Diesel engine aftertreatment systems, particularly the Selective Catalytic
Reduction (SCR) system with Ammonia Slip Catalyst (ASC), are critical for
reducing harmful $\nox$ emissions. The effectiveness of the SCR catalyst
degrades over time due to aging, leading to non-compliance with emission
standards. The existing approaches for on-board aging detection are often
impeded by the complexity of SCR dynamics, the cross-sensitivity of commercial
$\nox$ sensors to ammonia, and the general lack of on-board ammonia sensors in
commercial vehicles. This work presents a novel approach for modeling the SCR
system and estimating its parameters to diagnose catalyst aging, even with
uncertain and limited data from commercial trucks.

We propose a discrete-time, input-output model for $\nox$ reduction derived from
molar conservation across the control volume of SCR-ASC with focus on the
dynamics when the catalyst is saturated with ammonia. In this "saturated" mode,
the $\nox$ reduction, $\eta(k) = u_1(k-1) - x_1(k)$ where $u_1$ is the inlet
$\nox$ concentration at step $(k-1)$ and $x_1$ is the outlet $\nox$
concentration at step $(k)$, becomes independent of urea dosing. The $\nox$
reduction is modeled as an explicit function of inlet $\nox$ concentration
$\lr{u_1}$, the residence time $\lr{\tau}$, which is a function of the mass
flow rate, and the average concentration of adsorbed ammonia $\lr{\gamma}$,
which strongly depends on the urea-dosing rate. Under catalyst saturation,
$\gamma$ reaches its maximum temperature-dependent capacity, $\Gamma(T)$. Thus,
the $\nox$ reduction in the saturated mode, $\eta_{sat}$, is: $ \eta_{sat}
    (k+1) = u_1(k) \tau(k) k_{scr}(k) \Gamma (k) $, where $k_{scr}$ is the SCR rate
constant. This saturated model is parameterized using Chebyshev polynomials for
temperature dependence, resulting in a linear-in-parameters regression model: $
    \eta_{sat} (k+1) = \phi_{sat}^T(k) \; \theta_{sat} $. A crucial observation is
that the $\nox$ reduction under saturation, $\eta_{sat}(k)$, forms a tight
upper bound for the actual $\nox$ reduction, $\eta(k)$, at any given time. $
    \eta(k) \leq \eta_{sat}(k) \: \forall k $. This inequality allows us to
estimate the saturated model parameters, $\theta_{sat}$, from operational data
where the catalyst's saturation state is unknown.

The proposed methodology is validated using test-cell data, demonstrating
consistent parameter estimates across datasets and robustness to sensor
cross-sensitivity. Furthermore, we show a clear, statistically-significant
difference in parameters between new and aged catalysts using both test-cell
and real-world truck data. Using this distinction, a Wald-test-based aging
detector is developed. The method is shown to be robust to sensor
cross-sensitivity and to provide a reliable metric for detecting catalyst aging,
paving the way for more accurate and robust on-board aging diagnostic systems.